\Needspace{12\baselineskip}
\section{GPU Kernels: FlashAttention-2}
\label{sec:flash-attention}

\problem{pytorch\_attention}{PyTorch Attention Benchmarking}{2}

\partquestion{a}{Benchmark forward and backward attention for batch size 8, embedding dimensions $[16,32,64,128]$, and sequence lengths $[256,1024,4096,8192,16384]$. Record out-of-memory cases, account for memory in one small failing case, describe saved-for-backward memory scaling, and explain how to eliminate that cost.}
\deliverable{A results table with timings and memory calculations, plus a 1--2 paragraph discussion.}
\begin{table}[htbp]
\centering\small
\begin{tabular}{@{}rrrrll@{}}
\toprule
$d_{model}$ & Sequence length & Forward time & Backward time & Memory estimate & Status \\
\midrule
 & & & & & \\
 & & & & & \\
 & & & & & \\
\bottomrule
\end{tabular}
\caption{PyTorch attention benchmark (batch size 8). Add rows for each configuration.}
\label{tab:pytorch-attention}
\end{table}
\answer

\problem{torch\_compile}{Torch Compile}{2}

\partquestion{a}{Compare compiled and uncompiled attention with the configurations above.}
\deliverable{A table comparing forward and backward timings.}
\begin{table}[htbp]
\centering\small
\begin{tabular}{@{}rrrrrr@{}}
\toprule
$d_{model}$ & Sequence length & Eager fwd & Compiled fwd & Eager bwd & Compiled bwd \\
\midrule
 & & & & & \\
 & & & & & \\
\bottomrule
\end{tabular}
\caption{Compiled versus eager attention timings.}
\label{tab:compiled-attention}
\end{table}

\partquestion{b}{Compile the full Transformer and compare forward and full training-step performance.}
\deliverable{A table comparing the vanilla and compiled Transformer.}
\begin{table}[htbp]
\centering\small
\begin{tabular}{@{}llrrr@{}}
\toprule
Model size & Mode & Forward time & Forward + backward time & Full step time \\
\midrule
 & Vanilla & & & \\
 & Compiled & & & \\
\bottomrule
\end{tabular}
\caption{End-to-end Transformer timings with and without compilation.}
\label{tab:compiled-transformer}
\end{table}

\problem{flash\_forward}{FlashAttention-2 Forward Pass}{15}

\partquestion{a}{Implement a PyTorch-only \texttt{torch.autograd.Function} forward pass that returns output and log-sum-exp, saving the required tensors for backward. Causal masking is not required in this part.}
\deliverable{The PyTorch autograd implementation and the specified forward-pass validation result.}
\answer

\partquestion{b}{Implement the tiled Triton forward kernel and an autograd wrapper that calls it.}
\deliverable{The Triton autograd implementation and the specified validation result.}
\answer

\partquestion{c}{Add an optional causal-mask flag to the Triton implementation and save the flag for backward.}
\deliverable{The causal-mask implementation with the default remaining non-causal.}
\answer

\problem{flash\_backward}{FlashAttention-2 Backward Pass}{5}

Implement the backward pass using a PyTorch function compiled with \texttt{torch.compile}. Use $Q,K,V,O,dO,L$, compute the $D$ vector, and return $dQ,dK,dV$.

\deliverable{The backward implementation and the specified validation result.}
\answer

\problem{flash\_benchmarking}{FlashAttention-2 Benchmarking}{5}

Compare the partially Triton FlashAttention-2 implementation with regular PyTorch attention using \texttt{triton.testing.do\_bench}, batch size 1, causal masking, a single B200, sequence lengths from 128 to 65536, dimensions from 16 to 128, and BF16 and FP32.

\deliverable{A table comparing forward, backward, and end-to-end latency for both implementations.}
\begin{table}[htbp]
\centering\small
\begin{tabular}{@{}llrrrrrr@{}}
\toprule
Sequence & Dimension & Dtype & PyTorch fwd & FA2 fwd & PyTorch bwd & FA2 bwd & End-to-end \\
\midrule
 & & & & & & & \\
 & & & & & & & \\
 & & & & & & & \\
\bottomrule
\end{tabular}
\caption{FlashAttention-2 benchmark results. Add rows for each configuration and separate end-to-end values by implementation.}
\label{tab:flash-benchmark}
\end{table}
